\documentclass[10pt]{extarticle}
\usepackage[a4paper,margin=1.5cm]{geometry}
\usepackage[T1]{fontenc}
\usepackage{amsmath,amssymb}
\usepackage{graphicx}
\usepackage[export]{adjustbox}
\usepackage{multicol}
\usepackage{enumitem}
\usepackage{booktabs,array,tabularx}
\setlength{\columnsep}{0.7cm}
\setlength{\parindent}{0pt}
\setlength{\parskip}{3pt}
\pagestyle{empty}
\makeatletter
% figure registry (figdefs.tex is generated by the extract step)
\newcommand{\qfigdef}[4]{\@namedef{qf@file@#1}{#2}\@namedef{qf@w@#1}{#3}\@namedef{qf@wide@#1}{#4}}
\newcommand{\qf@check}[1]{\@ifundefined{qf@file@#1}{\PackageError{score2}{Unknown figure id #1}{}}{}}
\newcommand{\qf@img}[4][]{\includegraphics[width=#2\dimexpr\@nameuse{qf@w@#4}pt\relax,max width=#3,max height=0.45\textheight,#1]{\@nameuse{qf@file@#4}}}
% Full-width questions: finalize wraps every question that contains a wide figure in \qfullbegin...\qfullend,
% so the WHOLE question (stem, figures, options) is set across the page in reading order, then the two columns
% resume. The question's head - its number, stem and first figure (\qfullsplit marks the end of the head) -
% is kept on one page, so a number is never left at a page foot with its figure overleaf; the rest may break
% across pages normally, so no big gaps are left. Without a \qfullsplit, the whole question is kept together
% if it is under 45% of a page. Never write these yourself.
\newif\ifqf@full
\newif\ifqf@boxed
\newsavebox{\qf@box}
\newcommand{\qf@open}{\global\qf@fulltrue\par\global\qf@boxedtrue
  \setbox\qf@box\vbox\bgroup\hsize\textwidth\linewidth\textwidth}
\newcommand{\qf@emit}[1]{\par\egroup\global\qf@boxedfalse
  \ifdim\dimexpr\ht\qf@box+\dp\qf@box\relax<#1\textheight\noindent\usebox{\qf@box}\par\else\unvbox\qf@box\fi}
\newcommand{\qfullsplit}{\ifqf@boxed\qf@emit{0.85}\fi}
\newcommand{\qf@close}{\ifqf@boxed\qf@emit{0.45}\fi\global\qf@fullfalse}
\newcommand{\qfullbegin}{\end{multicols}\qf@open}
\newcommand{\qfullend}{\qf@close\begin{multicols}{2}}
\newcommand{\qfullstart}{\qf@open}
\newcommand{\qfulldone}{\qf@close}
% \qfig[scale]{id}: block figure, centred, directly where it is written (size set by the tool).
% Fallback if a wide figure is not inside a full-width question: break out of the columns in place - no floats.
\newcommand{\qfig}[2][1]{\qf@check{#2}%
  \ifqf@full\qf@blk{#1}{#2}\else
  \ifnum\@nameuse{qf@wide@#2}=1 \expandafter\@firstoftwo\else\expandafter\@secondoftwo\fi
  {\end{multicols}\begin{center}\qf@img{#1}{\textwidth}{#2}\end{center}\begin{multicols}{2}}%
  {\qf@blk{#1}{#2}}\fi}
\newcommand{\qf@blk}[2]{\par\smallskip{\centering\qf@img{#1}{\linewidth}{#2}\par}\smallskip}
% \qfigi[scale]{id}: inline figure (inside option lists / tables / minipages); never breaks columns.
\newcommand{\qfigi}[2][1]{\qf@check{#2}\qf@img[valign=c]{#1}{\linewidth}{#2}}
\makeatother
% structure macros - lint reads these to check completeness and figure placement
\newcommand{\papersec}[1]{\par\bigskip{\centering\Large\bfseries PAPER-#1\par}\medskip}
\newcommand{\partsec}[2]{\par\medskip\noindent{\large\bfseries PART-#1: #2}\par\smallskip}
\newcommand{\sect}[1]{\par\smallskip\noindent\textbf{#1}\par}
\newcommand{\qn}[1]{\par\smallskip\noindent\textbf{#1.}~}
\newcommand{\paraq}[2]{\par\smallskip\noindent\textbf{Paragraph for Questions #1 and #2}\par}
\newenvironment{qoptions}{\begin{itemize}[leftmargin=1.8em,itemsep=1pt,topsep=2pt,label={}]}{\end{itemize}}
\newcommand{\opt}[2]{\item[(#1)] #2}
% figure registry (generated by score2_tools.py finalize)
\qfigdef{p1_phy_s1_q02_1}{TOAS - Test 09 (30 May 2022)-img/p1_phy_s1_q02_1.png}{118.2}{0}
\qfigdef{p1_phy_s1_q03_B}{TOAS - Test 09 (30 May 2022)-img/p1_phy_s1_q03_B.png}{118.0}{0}
\qfigdef{p1_phy_s1_q04_1}{TOAS - Test 09 (30 May 2022)-img/p1_phy_s1_q04_1.png}{95.1}{0}
\qfigdef{p1_phy_s2_q05_1}{TOAS - Test 09 (30 May 2022)-img/p1_phy_s2_q05_1.png}{110.4}{0}
\qfigdef{p1_phy_s2_q06_1}{TOAS - Test 09 (30 May 2022)-img/p1_phy_s2_q06_1.png}{149.2}{0}
\qfigdef{p1_phy_s2_q07_1}{TOAS - Test 09 (30 May 2022)-img/p1_phy_s2_q07_1.png}{156.7}{0}
\qfigdef{p1_phy_s2_q08_1}{TOAS - Test 09 (30 May 2022)-img/p1_phy_s2_q08_1.png}{130.6}{0}
\qfigdef{p1_phy_s2_q09_1}{TOAS - Test 09 (30 May 2022)-img/p1_phy_s2_q09_1.png}{124.6}{0}
\qfigdef{p1_phy_s2_q10_1}{TOAS - Test 09 (30 May 2022)-img/p1_phy_s2_q10_1.png}{76.2}{0}
\qfigdef{p1_phy_s3_para1_2_1}{TOAS - Test 09 (30 May 2022)-img/p1_phy_s3_para1_2_1.png}{98.4}{0}
\qfigdef{p1_phy_s3_para3_4_1}{TOAS - Test 09 (30 May 2022)-img/p1_phy_s3_para3_4_1.png}{95.5}{0}
\qfigdef{p1_phy_s3_para5_6_1}{TOAS - Test 09 (30 May 2022)-img/p1_phy_s3_para5_6_1.png}{127.3}{0}
\qfigdef{p1_phy_s4_q07_1}{TOAS - Test 09 (30 May 2022)-img/p1_phy_s4_q07_1.png}{79.1}{0}
\qfigdef{p1_phy_s4_q09_1}{TOAS - Test 09 (30 May 2022)-img/p1_phy_s4_q09_1.png}{337.2}{1}
\qfigdef{p1_chem_s1_q01_A}{TOAS - Test 09 (30 May 2022)-img/p1_chem_s1_q01_A.png}{75.4}{0}
\qfigdef{p1_chem_s1_q01_B}{TOAS - Test 09 (30 May 2022)-img/p1_chem_s1_q01_B.png}{91.1}{0}
\qfigdef{p1_chem_s1_q01_C}{TOAS - Test 09 (30 May 2022)-img/p1_chem_s1_q01_C.png}{84.0}{0}
\qfigdef{p1_chem_s1_q01_D}{TOAS - Test 09 (30 May 2022)-img/p1_chem_s1_q01_D.png}{90.1}{0}
\qfigdef{p1_chem_s1_q02_1}{TOAS - Test 09 (30 May 2022)-img/p1_chem_s1_q02_1.png}{108.7}{0}
\qfigdef{p1_chem_s1_q02_A}{TOAS - Test 09 (30 May 2022)-img/p1_chem_s1_q02_A.png}{62.5}{0}
\qfigdef{p1_chem_s1_q02_B}{TOAS - Test 09 (30 May 2022)-img/p1_chem_s1_q02_B.png}{62.5}{0}
\qfigdef{p1_chem_s1_q02_C}{TOAS - Test 09 (30 May 2022)-img/p1_chem_s1_q02_C.png}{60.4}{0}
\qfigdef{p1_chem_s1_q02_D}{TOAS - Test 09 (30 May 2022)-img/p1_chem_s1_q02_D.png}{47.3}{0}
\qfigdef{p1_chem_s2_q08_A}{TOAS - Test 09 (30 May 2022)-img/p1_chem_s2_q08_A.png}{35.8}{0}
\qfigdef{p1_chem_s2_q08_B}{TOAS - Test 09 (30 May 2022)-img/p1_chem_s2_q08_B.png}{48.7}{0}
\qfigdef{p1_chem_s2_q08_C}{TOAS - Test 09 (30 May 2022)-img/p1_chem_s2_q08_C.png}{37.7}{0}
\qfigdef{p1_chem_s2_q08_D}{TOAS - Test 09 (30 May 2022)-img/p1_chem_s2_q08_D.png}{28.2}{0}
\qfigdef{p1_chem_s2_q09_1}{TOAS - Test 09 (30 May 2022)-img/p1_chem_s2_q09_1.png}{183.4}{0}
\qfigdef{p1_chem_s3_para5_6_1}{TOAS - Test 09 (30 May 2022)-img/p1_chem_s3_para5_6_1.png}{318.2}{1}
\qfigdef{p1_chem_s4_q09_1}{TOAS - Test 09 (30 May 2022)-img/p1_chem_s4_q09_1.png}{166.8}{0}
\qfigdef{p2_phy_s1_q01_1}{TOAS - Test 09 (30 May 2022)-img/p2_phy_s1_q01_1.png}{151.8}{0}
\qfigdef{p2_phy_s1_q02_1}{TOAS - Test 09 (30 May 2022)-img/p2_phy_s1_q02_1.png}{142.1}{0}
\qfigdef{p2_phy_s1_q04_1}{TOAS - Test 09 (30 May 2022)-img/p2_phy_s1_q04_1.png}{129.8}{0}
\qfigdef{p2_phy_s1_q05_1}{TOAS - Test 09 (30 May 2022)-img/p2_phy_s1_q05_1.png}{151.8}{0}
\qfigdef{p2_phy_s1_q06_1}{TOAS - Test 09 (30 May 2022)-img/p2_phy_s1_q06_1.png}{152.0}{0}
\qfigdef{p2_phy_s1_q07_1}{TOAS - Test 09 (30 May 2022)-img/p2_phy_s1_q07_1.png}{127.3}{0}
\qfigdef{p2_phy_s1_q08_1}{TOAS - Test 09 (30 May 2022)-img/p2_phy_s1_q08_1.png}{92.0}{0}
\qfigdef{p2_phy_s1_q09_1}{TOAS - Test 09 (30 May 2022)-img/p2_phy_s1_q09_1.png}{59.6}{0}
\qfigdef{p2_phy_s1_q10_1}{TOAS - Test 09 (30 May 2022)-img/p2_phy_s1_q10_1.png}{112.0}{0}
\qfigdef{p2_phy_s1_q11_1}{TOAS - Test 09 (30 May 2022)-img/p2_phy_s1_q11_1.png}{101.1}{0}
\qfigdef{p2_phy_s1_q12_1}{TOAS - Test 09 (30 May 2022)-img/p2_phy_s1_q12_1.png}{150.6}{0}
\qfigdef{p2_phy_s1_q13_1}{TOAS - Test 09 (30 May 2022)-img/p2_phy_s1_q13_1.png}{100.9}{0}
\qfigdef{p2_phy_s1_q14_1}{TOAS - Test 09 (30 May 2022)-img/p2_phy_s1_q14_1.png}{103.5}{0}
\qfigdef{p2_phy_s1_q15_A}{TOAS - Test 09 (30 May 2022)-img/p2_phy_s1_q15_A.png}{77.7}{0}
\qfigdef{p2_phy_s1_q16_1}{TOAS - Test 09 (30 May 2022)-img/p2_phy_s1_q16_1.png}{109.5}{0}
\qfigdef{p2_phy_s1_q18_1}{TOAS - Test 09 (30 May 2022)-img/p2_phy_s1_q18_1.png}{139.1}{0}
\qfigdef{p2_phy_s1_q20_1}{TOAS - Test 09 (30 May 2022)-img/p2_phy_s1_q20_1.png}{110.0}{0}
\qfigdef{p2_phy_s1_q20_A}{TOAS - Test 09 (30 May 2022)-img/p2_phy_s1_q20_A.png}{76.1}{0}
\qfigdef{p2_phy_s1_q20_B}{TOAS - Test 09 (30 May 2022)-img/p2_phy_s1_q20_B.png}{76.2}{0}
\qfigdef{p2_phy_s1_q20_C}{TOAS - Test 09 (30 May 2022)-img/p2_phy_s1_q20_C.png}{75.6}{0}
\qfigdef{p2_phy_s1_q20_D}{TOAS - Test 09 (30 May 2022)-img/p2_phy_s1_q20_D.png}{79.8}{0}
\qfigdef{p2_phy_s2_q02_1}{TOAS - Test 09 (30 May 2022)-img/p2_phy_s2_q02_1.png}{60.7}{0}
\qfigdef{p2_phy_s2_q03_1}{TOAS - Test 09 (30 May 2022)-img/p2_phy_s2_q03_1.png}{99.7}{0}
\qfigdef{p2_phy_s2_q04_1}{TOAS - Test 09 (30 May 2022)-img/p2_phy_s2_q04_1.png}{84.9}{0}
\qfigdef{p2_phy_s2_q05_1}{TOAS - Test 09 (30 May 2022)-img/p2_phy_s2_q05_1.png}{113.4}{0}
\qfigdef{p2_phy_s2_q06_1}{TOAS - Test 09 (30 May 2022)-img/p2_phy_s2_q06_1.png}{128.3}{0}
\qfigdef{p2_phy_s2_q07_1}{TOAS - Test 09 (30 May 2022)-img/p2_phy_s2_q07_1.png}{134.5}{0}
\qfigdef{p2_phy_s2_q08_1}{TOAS - Test 09 (30 May 2022)-img/p2_phy_s2_q08_1.png}{123.9}{0}
\qfigdef{p2_phy_s2_q09_1}{TOAS - Test 09 (30 May 2022)-img/p2_phy_s2_q09_1.png}{153.7}{0}
\qfigdef{p2_phy_s2_q10_1}{TOAS - Test 09 (30 May 2022)-img/p2_phy_s2_q10_1.png}{106.5}{0}
\qfigdef{p2_chem_s1_q07_A}{TOAS - Test 09 (30 May 2022)-img/p2_chem_s1_q07_A.png}{91.3}{0}
\qfigdef{p2_chem_s1_q07_C}{TOAS - Test 09 (30 May 2022)-img/p2_chem_s1_q07_C.png}{101.5}{0}
\qfigdef{p2_chem_s1_q07_D}{TOAS - Test 09 (30 May 2022)-img/p2_chem_s1_q07_D.png}{102.7}{0}
\qfigdef{p2_chem_s1_q08_A}{TOAS - Test 09 (30 May 2022)-img/p2_chem_s1_q08_A.png}{93.8}{0}
\qfigdef{p2_chem_s1_q08_B}{TOAS - Test 09 (30 May 2022)-img/p2_chem_s1_q08_B.png}{125.2}{0}
\qfigdef{p2_chem_s1_q08_C}{TOAS - Test 09 (30 May 2022)-img/p2_chem_s1_q08_C.png}{98.0}{0}
\qfigdef{p2_chem_s1_q08_4}{TOAS - Test 09 (30 May 2022)-img/p2_chem_s1_q08_4.png}{129.8}{0}
\qfigdef{p2_chem_s1_q09_1}{TOAS - Test 09 (30 May 2022)-img/p2_chem_s1_q09_1.png}{145.8}{0}
\qfigdef{p2_chem_s1_q09_A}{TOAS - Test 09 (30 May 2022)-img/p2_chem_s1_q09_A.png}{21.6}{0}
\qfigdef{p2_chem_s1_q11_A}{TOAS - Test 09 (30 May 2022)-img/p2_chem_s1_q11_A.png}{84.6}{0}
\qfigdef{p2_chem_s1_q11_B}{TOAS - Test 09 (30 May 2022)-img/p2_chem_s1_q11_B.png}{48.3}{0}
\qfigdef{p2_chem_s1_q11_C}{TOAS - Test 09 (30 May 2022)-img/p2_chem_s1_q11_C.png}{49.5}{0}
\qfigdef{p2_chem_s1_q11_D}{TOAS - Test 09 (30 May 2022)-img/p2_chem_s1_q11_D.png}{52.0}{0}
\qfigdef{p2_chem_s1_q12_1}{TOAS - Test 09 (30 May 2022)-img/p2_chem_s1_q12_1.png}{109.2}{0}
\qfigdef{p2_chem_s1_q12_A}{TOAS - Test 09 (30 May 2022)-img/p2_chem_s1_q12_A.png}{60.0}{0}
\qfigdef{p2_chem_s1_q12_B}{TOAS - Test 09 (30 May 2022)-img/p2_chem_s1_q12_B.png}{54.6}{0}
\qfigdef{p2_chem_s1_q12_C}{TOAS - Test 09 (30 May 2022)-img/p2_chem_s1_q12_C.png}{48.5}{0}
\qfigdef{p2_chem_s1_q12_D}{TOAS - Test 09 (30 May 2022)-img/p2_chem_s1_q12_D.png}{58.0}{0}
\qfigdef{p2_chem_s1_q13_B}{TOAS - Test 09 (30 May 2022)-img/p2_chem_s1_q13_B.png}{69.3}{0}
\qfigdef{p2_chem_s2_q05_1}{TOAS - Test 09 (30 May 2022)-img/p2_chem_s2_q05_1.png}{77.9}{0}
\qfigdef{p2_chem_s2_q06_1}{TOAS - Test 09 (30 May 2022)-img/p2_chem_s2_q06_1.png}{134.0}{0}
\qfigdef{p2_chem_s2_q07_1}{TOAS - Test 09 (30 May 2022)-img/p2_chem_s2_q07_1.png}{149.4}{0}
\raggedcolumns
\makeatletter\@beginparpenalty=10000\makeatother
\newcolumntype{L}[1]{>{\raggedright\arraybackslash}p{#1}}
\begin{document}
\begin{center}{\Large\bfseries TOAS --- Test 09 (30 May 2022)}\end{center}
\begin{multicols}{2}
\raggedcolumns
\papersec{1}
\partsec{1}{PHYSICS}
\sect{SECTION-I(i)}
\qn{1} Moment of inertia of a uniform equilateral triangular lamina of mass 20 kg and side length 75 cm about an axis perpendicular to plane and passing through centroid (in kg--m$^2$) is :-
\begin{qoptions}
\opt{A}{$\dfrac{360}{96}$}
\opt{B}{$\dfrac{360}{32}$}
\opt{C}{$\dfrac{360}{192}$}
\opt{D}{$\dfrac{180}{192}$}
\end{qoptions}
\qn{2} A battery of emf 48 V unknown internal resistance is connected to a rheostat whose resistance can be varied in a long range. A graph is plotted between power developed across rheostat and resistance of rheostat which is shown in figure. The internal resistance of the battery in $\Omega$ is :-
\qfig{p1_phy_s1_q02_1}
\begin{qoptions}
\opt{A}{1.44}
\opt{B}{2.88}
\opt{C}{5.76}
\opt{D}{11.52}
\end{qoptions}
\qn{3} All the springs and pulleys are ideal and massless. If the block is slightly pulled vertically downward from its equilibrium position, then the time period of oscillation of the block is (assume strings are taught all the time)
\qfig{p1_phy_s1_q03_B}
\begin{qoptions}
\opt{A}{$2\pi\sqrt{\dfrac{m}{4k}}$}
\opt{B}{$\pi\sqrt{\dfrac{2m}{k}}$}
\opt{C}{$\pi\sqrt{\dfrac{m}{2k}}$}
\opt{D}{$2\pi\sqrt{\dfrac{m}{16k}}$}
\end{qoptions}
\qn{4} A potentiometer wire AB is 100 cm long and has total resistance of 10$\Omega$. Find the value of unknown resistance R so that null point is obtained at a distance 40 cm from A.
\qfig{p1_phy_s1_q04_1}
\begin{qoptions}
\opt{A}{1$\Omega$}
\opt{B}{2$\Omega$}
\opt{C}{3$\Omega$}
\opt{D}{4$\Omega$}
\end{qoptions}
\sect{SECTION-I(ii)}
\qn{5} Two surface OABC and OCDE lies in the plane of xy and yz as shown in the figure. A charged particle `q' lies in the space at a point P, if:
\qfig{p1_phy_s2_q05_1}
\begin{qoptions}
\opt{A}{coordinates of `P' is $(a-\Delta r,\ a-\Delta r,\ \Delta r)$ and $a\gg\Delta r$, then flux passing through surface OABC is $\dfrac{7q}{24\varepsilon_0}$.}
\opt{B}{coordinates of `P' is $(a-\Delta r,\ a-\Delta r,\ \Delta r)$ and $a\gg\Delta r$, then flux passing through surface OCDE is $\dfrac{q}{24\varepsilon_0}$.}
\opt{C}{coordinates of `P' is $(a+\Delta r,\ a+\Delta r,\ -\Delta r)$ and $a\gg\Delta r$, then flux passing through surface OABC is $\dfrac{q}{24\varepsilon_0}$.}
\opt{D}{coordinates of `P' is $(a+\Delta r,\ a+\Delta r,\ -\Delta r)$ and $a\gg\Delta r$, then flux passing through surface OCDE is $\dfrac{q}{24\varepsilon_0}$.}
\end{qoptions}
\qn{6} The resistance of wire PQ is 410 $\Omega$, which is tapped with the external wire as shown in the figure. The value of R is not known. If the potential difference between AB is 400 volts and the net current through the source battery is 6 ampere, then which of the following is/are correct :-
\qfig{p1_phy_s2_q06_1}
\begin{qoptions}
\opt{A}{The resistance of unknown resistor R is 200 $\Omega$.}
\opt{B}{The equivalent resistance across AB is $\dfrac{200}{3}\,\Omega$.}
\opt{C}{The current in 25$\Omega$ resistor is $\dfrac{2}{5}$ ampere.}
\opt{D}{The potential drop across 120$\Omega$ resistor is 60 volts.}
\end{qoptions}
\qn{7} A non conducting uniform thin shell of mass $m$ and radius $R$ is having uniform charge density $+\sigma$ on one half and $-\sigma$ on another half as shown in figure. It is placed on a rough non conducting horizontal plane. At $t=0$ a uniform electric field $\vec E=-E_0\hat j$ N/C is switched on and the solid sphere starts rolling without slipping. Then which of the following is/are correct.
\qfig{p1_phy_s2_q07_1}
\begin{qoptions}
\opt{A}{The acceleration of the topmost point of the hollow sphere at $t=0$, is $\dfrac{36}{25}\dfrac{\pi\sigma E_0R^2}{m}$.}
\opt{B}{The magnitude of frictional force at $t=0$ is $\dfrac{18}{25}\pi\sigma E_0R^2$.}
\opt{C}{The acceleration of centre of mass of the hollow sphere at $t=0$, is $\dfrac{12}{25}\dfrac{\pi\sigma E_0R^2}{m}$.}
\opt{D}{The magnitude of frictional force at $t=0$ is $\dfrac{9}{25}\pi\sigma E_0R^2$.}
\end{qoptions}
\qn{8} Consider a satellite `S' which is just above the ECUADOR and rotating on equatorial plane of earth. The sense of rotation of earth and satellite are similar means both are rotating from west to east. The time period of satellite is 6 hours. At $t=0$, satellite is just above the Brazil. Now choose the correct option (s).
\qfig{p1_phy_s2_q08_1}
\begin{qoptions}
\opt{A}{The time after which the satellite will be above the ECUADOR again is 8 hours.}
\opt{B}{The above satellite may be a geostationary satellite.}
\opt{C}{The height of geostationary satellite above the sea level is 35,800 km approximately.}
\opt{D}{If the distance from the centre of earth to the satellite `S' becomes $4r_0$, then the time after which the satellite will be above the ECUADOR again is 24 hours.}
\end{qoptions}
\qn{9} A tunnel is dug across the earth along a chord AB as shown, A particle is projected inside the tunnel from end `A' with an initial velocity $v_0=\dfrac{\sqrt{3g_0R}}{2}$, there is no friction between the tunnel and the particle. Then choose the correct option (s), (where $g_0$ = gravitational acceleration on the surface of Earth, $R$ = radius of earth)
\qfig{p1_phy_s2_q09_1}
\begin{qoptions}
\opt{A}{The time taken by the particle to move from end `A' to other end `B' of the tunnel is $\dfrac{\pi}{3}\sqrt{\dfrac{R}{g_0}}$.}
\opt{B}{The time taken by the particle to move from end `A' to other end `B' of the tunnel is $\dfrac{2\pi}{3}\sqrt{\dfrac{R}{g_0}}$.}
\opt{C}{The maximum speed of the particle inside the tunnel is $\sqrt{g_0R}$.}
\opt{D}{The maximum speed of the particle inside the tunnel is $\sqrt{2g_0R}$.}
\end{qoptions}
\qn{10} A bird is flying downward with constant velocity 6 cm/s towards a fixed container filled with a liquid of refractive index $\dfrac43$. A fish is moving upward with constant speed 4 cm/s towards the surface of liquid. Surface of liquid is also moving downward with constant velocity 2 cm/s shown in the figure. The bottom of container is silvered. Then :-
\qfig{p1_phy_s2_q10_1}
\begin{qoptions}
\opt{A}{Speed of fish as seen by bird only after refraction is 9 cm/sec.}
\opt{B}{Speed of fish as seen by bird only after refraction is 8.5 cm/sec.}
\opt{C}{Speed of fish as seen by bird after reflection and then refraction is 3 cm/sec.}
\opt{D}{Speed of fish as seen by bird after reflection and then refraction is 2.5 cm/sec.}
\end{qoptions}
\sect{SECTION-II(i)}
\paraq{1}{2}
A uniform circular disc of mass $M=4$ kg and radius $R=25$ cm is suspended in a vertical plane and hinged at its centre `O'. It is free to rotate about a horizontal axis passing through its centre and perpendicular to its plane. A particle of mass $m=0.5$ kg moving horizontally with a velocity $u=20$ m/s collides with disc at point `P' and gets stuck to it as shown. ($g=10$ m/s$^2$)
\qfig{p1_phy_s3_para1_2_1}
\qn{1} Find the angular velocity of the disc just after collision (in rad/sec)
\qn{2} The force on the disc due to hinge at its centre when the particle reaches the highest point (in newton).
\paraq{3}{4}
A converging lens of power $P_1=+5$D and a concave mirror of power $P_2=+16$D are placed at a separation 80 cm such that their principal axis coincide. A point object is placed left side of the lens on its principal axes at a distance $X_1$ from the lens such that image of the object coincides with the object itself. When the mirror is further shifted right side by a distance $X_2$, the image of the object further coincides with the object.
\qfig{p1_phy_s3_para3_4_1}
\qn{3} The value of $X_1$ (in cm) is :
\qn{4} The value of $X_2$ (in cm) is:
\paraq{5}{6}
In the circuit shown in figure, $R_1=R_2=R_3=3\,\Omega$ and e.m.f. of each cell is $E=4$V and negligible internal resistance. All ammeters are ideal.
\qfig{p1_phy_s3_para5_6_1}
\qn{5} Reading of ammeter $A_1$ in ampere is:
\qn{6} Potential difference (in Volt) between points a and b is
\sect{SECTION-II(ii)}
\qn{7} A small body of mass `m' is projected from the surface of a given planet as shown in the figure. If small body has velocity which is just sufficient to get out from the gravity of planet then time taken by small body to reach at height of $3R$ from the surface of planet is $\dfrac{K}{3}\sqrt{\dfrac{2R^3}{GM}}$. Find the value of K.
\qfig{p1_phy_s4_q07_1}
\qn{8} A non-conducing solid sphere has a charge density given by $\rho=\rho_0r^n$, where $\rho_0$, $n$ are constant and `r' is the distance from the centre of sphere. If the magnitude of electric field at half of the radius is $\dfrac{1}{16}$ times that of the surfaces of nonconducting then find the value of $n$.
\qfullbegin
\qn{9} Find the velocity (in cm/s) of the car C as appeared to the observer shown in the figure. There are several parallel slabs of different medium placed between the car and the observer.
\qfig{p1_phy_s4_q09_1}
\qfullsplit
\raggedcolumns
\qfullend
\partsec{2}{CHEMISTRY}
\sect{SECTION-I(i)}
\qn{1} Which of the following will give non resolvable (meso) product :
\begin{qoptions}
\opt{A}{\qfigi{p1_chem_s1_q01_A}}
\opt{B}{\qfigi{p1_chem_s1_q01_B}}
\opt{C}{\qfigi{p1_chem_s1_q01_C}}
\opt{D}{\qfigi{p1_chem_s1_q01_D}}
\end{qoptions}
\qn{2} Which compound is suitable for `P' in following reaction :
\par\smallskip P $\xrightarrow[\mathrm{H_2SO_4}]{\mathrm{HgSO_4}}\ \xrightarrow[\mathrm{H_2O}]{\mathrm{NaBH_4}}$ \qfigi{p1_chem_s1_q02_1}
\begin{qoptions}
\opt{A}{\qfigi{p1_chem_s1_q02_A}}
\opt{B}{\qfigi{p1_chem_s1_q02_B}}
\opt{C}{\qfigi{p1_chem_s1_q02_C}}
\opt{D}{\qfigi{p1_chem_s1_q02_D}}
\end{qoptions}
\qn{3} Select the incorrect statement :
\begin{qoptions}
\opt{A}{MgFe$_2$O$_4$ can behave like a ferrimagnetic solid}
\opt{B}{The unit cell having crystal parameters, $a=b\neq c$, $\alpha=\beta=90^\circ$, $\gamma=120^\circ$ is hexagonal}
\opt{C}{In ionic compounds having Frenkel defect, the ratio $r^+/r^-$ is high}
\opt{D}{The co-ordination number of Na$^+$ ion in NaCl is 6}
\end{qoptions}
\qn{4} If FeO crystallizes rock salt structure in which O$^{2-}$ ions occupy FCC lattice points. Due to some crystal defect, lattice becomes Fe$_{0.75}$O. Find number of effective Fe$^{3+}$ ions per unit cell -
\begin{qoptions}
\opt{A}{1}
\opt{B}{2}
\opt{C}{3}
\opt{D}{4}
\end{qoptions}
\sect{SECTION-I(ii)}
\qn{5} The CORRECT order(s) regarding the atomic size is / are
\begin{qoptions}
\opt{A}{Zr $\simeq$ Hf}
\opt{B}{Fe $\simeq$ Co $\simeq$ Ni}
\opt{C}{Sc $<$ Y $<$ La}
\opt{D}{Cu $<$ Zn}
\end{qoptions}
\qn{6} Which of the following compound(s) give(s) dibasic acid as a product on hydrolysis.
\begin{qoptions}
\opt{A}{P$_4$O$_6$}
\opt{B}{SO$_2$Cl$_2$}
\opt{C}{SF$_6$}
\opt{D}{SeF$_6$}
\end{qoptions}
\qn{7} Which of the following ore(s) contain(s) copper metal?
\begin{qoptions}
\opt{A}{Dolomite}
\opt{B}{Chalcopyrite}
\opt{C}{Fool's gold}
\opt{D}{Malachite}
\end{qoptions}
\qn{8} Which of the following will give same product with H$^{\oplus}$, KMnO$_4$, $\Delta$ and O$_3$ / Zn, H$_2$O ?
\begin{qoptions}
\opt{A}{\qfigi{p1_chem_s2_q08_A}}
\opt{B}{\qfigi{p1_chem_s2_q08_B}}
\opt{C}{\qfigi{p1_chem_s2_q08_C}}
\opt{D}{\qfigi{p1_chem_s2_q08_D}}
\end{qoptions}
\qn{9} \qfig{p1_chem_s2_q09_1}
Which of the following statement is correct.\par
(i) If A \& B are same gases then T$_2>$ T$_1$\par
(ii) If T$_1$ = T$_2$ then A$_{(g)}$ may be SO$_{2(g)}$ \& B$_{(g)}$ may be CH$_{4(g)}$\par
(iii) If B \& C are same gases then T$_3>$ T$_2$\par
(iv) If T$_1$ = T$_2$ = T$_3$ then A$_{(g)}$ may be CH$_4$, B$_{(g)}$ may be SO$_2$ \& C$_{(g)}$ may be SO$_{3(g)}$
\begin{qoptions}
\opt{A}{i, ii \& iii only}
\opt{B}{i, ii, \& iv only}
\opt{C}{All are correct}
\opt{D}{ii \& iii only}
\end{qoptions}
\qn{10} Two volatile liquids A \& B are mixed to form an ideal solution. If $P_A^0=100$ mm of Hg \& $P_B^0=150$ mm of Hg then identify the correct option(s) if 10 moles of A are mixed with 40 moles of B such that moles vaporised are very small.
\begin{qoptions}
\opt{A}{Total vapour pressure will be 140 mm of Hg.}
\opt{B}{Mole fraction of A in vapour state will be less than 0.5.}
\opt{C}{Mole fraction of B in vapour state is equal to $\dfrac{6}{7}$}
\opt{D}{Mole \% of B in vapour state will be greater than 80\%.}
\end{qoptions}
\sect{SECTION-II(i)}
\paraq{1}{2}
The complex compounds show space isomerism when the ligand species have different spatial arrangement surrounding the CMA
\qn{1} What is the number of stereo isomeric forms for the complex [Pt (py) (NH$_3$) (Br) (I)]
\qn{2} How many stereo isomers are possible for [Cr(H$_2$O)$_2$ (C$_2$O$_4$)$_2$]$^{-}$ ?
\paraq{3}{4}
Kinetics of acidic hydrolysis of ester is pseudo I order reaction
\[\begin{aligned}&\mathrm{CH_3COOCH_3(aq.)+\underset{excess}{H_2O}\,(\ell)}\\&\qquad\xrightarrow{[H^+]}\mathrm{CH_3COOH(aq.)+CH_3OH(aq.)}\end{aligned}\]
Rate law is given by : $r=K[\mathrm{H^+}][\mathrm{H_2O}][\mathrm{CH_3COOCH_3}]$ ; $K=1.8\times10^{-3}\ \mathrm{M^{-2}\,sec^{-1}}$ where H$^+$ ion concentration is given by acid catalyst. In an experiment 1M CH$_3$COOCH$_3$ is hydrolysed using 0.1M HCl as catalyst. [Use molarity of H$_2$O = $\dfrac{1000}{18}$M]
\qn{3} Reaction mixture is titrated against standard NaOH solution at different times. If initially 10 ml of NaOH was required while after long time 100 ml of NaOH is required then find volume of NaOH (in ml) required at 230.3 sec.
\qn{4} Calculate the time at which ester concentration reduce to 0.25M (in seconds) [ln2 = 0.7]
\qfullbegin
\paraq{5}{6}
Among the following alkenes :
\qfig{p1_chem_s3_para5_6_1}
\qfullend
\qn{5} How many alkenes which give meso compound on ozonolysis in presence of Zn?
\qn{6} From the following options identify alkenes which give optically active compound as one of the product on ozonolysis in presence of Zn.
\sect{SECTION-II(ii)}
\qn{7} Total number of complex compound which have low spin as well as Diamagnetic.\par
[Fe(CN)$_6$]$^{3-}$, [Co(H$_2$O)$_6$]$^{3\oplus}$, [NiF$_6$]$^{2-}$, [Cr(H$_2$O)$_6$]$^{3\oplus}$, [Fe(CN)$_6$]$^{4-}$, [Co(C$_2$O$_4$)$_3$]$^{3-}$, [Ni(en)$_3$]$^{3+}$, [PtCl$_4$]$^{2-}$
\qn{8} Total number of species in which highest occupied molecular orbital has gerade type symmetry.\par
B$_2$, N$_2$, C$_2$, N$_2^{2+}$, O$_2$, O$_2^{2-}$, N$_2^{2-}$
\qn{9} How many of the following will give silver mirror with Tollen's reagent ?\par
(i) Ph--CH=O, (ii) HCHO, (v) HCOOH, (vi) MeCOOH, (viii) HC$\equiv$CH, and the structures below, in order (iii), (iv), (vii), (ix), (x):
\qfig{p1_chem_s4_q09_1}
\raggedcolumns
\partsec{3}{MATHEMATICS}
\sect{SECTION-I(i)}
\qn{1} {\small\[\begin{aligned}\int\Big(&2e^{\tan^{-1}x+x^2+\ln x}\\&+\frac{2x^3+2x+1}{x+\frac1x}\,e^{\tan^{-1}x+x^2+\ln x}\Big)dx\end{aligned}\]}
is equal to (where C is constant of integration)
\begin{qoptions}
\opt{A}{$\dfrac{x^3}{3}e^{\tan^{-1}x+x^2+\ln x}+C$}
\opt{B}{$\dfrac{x^2}{2}e^{\tan^{-1}x+x^2+\ln x}+C$}
\opt{C}{$xe^{\tan^{-1}x+x^3+\ln x}+C$}
\opt{D}{$x^2e^{\tan^{-1}x+x^2}+C$}
\end{qoptions}
\qn{2} The function $f(x)=\alpha x^2-\beta x-4x^3+\gamma$, where $\alpha,\beta,\gamma\in R$, has local minima at $P(\log_2a,\ f(\log_2a))$ \& local maxima at $Q(\log_2a^2,\ f(\log_2a^2))$. If the graph of $f(x)$ changes concavity about the point $R\!\left(\dfrac34,\ f\!\left(\dfrac34\right)\right)$, then which of the following conic section can have eccentricity `a' -
\begin{qoptions}
\opt{A}{circle}
\opt{B}{parabola}
\opt{C}{ellipse}
\opt{D}{hyperbola}
\end{qoptions}
\qn{3} If $m$ is the slope of a line which is tangent to $y^3=x^4$ \& a normal to $x^2-2x+y^2=0$, then $\left(\dfrac{3m}{4}\right)^3$ is equal to ($m\neq0$)
\begin{qoptions}
\opt{A}{3}
\opt{B}{$\dfrac43$}
\opt{C}{4}
\opt{D}{$\dfrac34$}
\end{qoptions}
\qn{4} For the equation $b\sqrt{1-a^2}-a\sqrt{1-b^2}=\sin\!\left(\dfrac{b-a}{\sqrt{1-x^2}}\right)$, where $0<a<b<1/2$, then which of the following is true -
\begin{qoptions}
\opt{A}{there is no real root of the equation in $(a,b)$}
\opt{B}{there is exactly one real root of the equation in $(a,b)$}
\opt{C}{there is atleast one root of the equation in $(a,b)$}
\opt{D}{none of these}
\end{qoptions}
\sect{SECTION-I(ii)}
\qn{5} Consider twice derivable bijective function defined on $[0,1]$ such that $f(\alpha)=\dfrac13$ and $f(f(\beta))=\dfrac13$. If $f''(x)>0\ \forall x\in\left[0,\dfrac12\right)$ and $f''(x)<0\ \forall x\in\left(\dfrac12,1\right]$ with $f\!\left(\dfrac12\right)=\dfrac12$. Which the following holds good ?
\begin{qoptions}
\opt{A}{$\beta>\alpha$ implies $f(x)$ is an increasing function}
\opt{B}{$\beta>\alpha$ implies $f(x)$ is a decreasing function}
\opt{C}{$2f\!\left(\dfrac14\right)<f(0)+f\!\left(\dfrac12\right)$}
\opt{D}{$2f\!\left(\dfrac34\right)<f\!\left(\dfrac12\right)+f(1)$}
\end{qoptions}
\qn{6} For the function $f(x)=e^{e^x}-\tan^{-1}x-3$ which of the following holds good ?
\begin{qoptions}
\opt{A}{$f(x)=0$ has exactly one positive real root}
\opt{B}{$y=f(x)$ has exactly one minima at negative $x$}
\opt{C}{$y=f(x)$ has exactly one point of minima at positive $x$}
\opt{D}{$y=f(x)$ is strictly increasing function}
\end{qoptions}
\qn{7} Let $x_1,x_2,x_3$ be the roots of the equation $x(3-x^2)=1$, then
\begin{qoptions}
\opt{A}{exactly one of $x_1,x_2,x_3$ is real}
\opt{B}{exactly one of $x_1,x_2,x_3$ is rational}
\opt{C}{the value of $\displaystyle\sum_{i=1}^{3}\left(\lim_{m\to\infty}\left(\lim_{n\to\infty}\cos(\pi n!\,x_i)\right)^{m}\right)$ where $m,n\in N$, is 3}
\opt{D}{the value of $\displaystyle\sum_{i=1}^{3}\left(\lim_{m\to\infty}\left(\lim_{n\to\infty}\cos(\pi n!\,x_i)\right)^{m}\right)$ where $m,n\in N$, is 0}
\end{qoptions}
\qn{8} Let $f(x)$ be a non-negative continuous and bounded function for all $x\geqslant0$. If $(\cos x)f'(x)\leqslant(\sin x-\cos x)f(x)$ ; $\forall x\geqslant0$, then which of the following is / are correct ?
\begin{qoptions}
\opt{A}{$f(6)+f(5)>0$}
\opt{B}{$x^2-3x+2+f(7)=0$ has 2 distinct solutions}
\opt{C}{$f(5)f(7)-f(6)f(5)=0$}
\opt{D}{$\displaystyle\lim_{x\to6}\dfrac{f(x)-\sin(\pi x)}{(x-6)}=1$}
\end{qoptions}
\qn{9} Let $g(x)$ be a polynomial of degree 2018 and
\[f(x)=\begin{cases}g\!\left(\dfrac{1}{x^5}\right)e^{-\frac{1}{x^6}}, & x\neq0\\[4pt] 0, & x=0\end{cases}\]
then
\begin{qoptions}
\opt{A}{$f(x)$ is continuous at $x=0$}
\opt{B}{$f(x)$ is non-differentaible at $x=0$}
\opt{C}{$f'(0)=\displaystyle\lim_{x\to0}f(x)$}
\opt{D}{$\displaystyle\lim_{x\to0}f(x)$ does not exist}
\end{qoptions}
\qn{10} Let
\[f(x)=\begin{cases}-2 & \text{for } -3\leqslant x\leqslant0\\ x-2 & \text{for } 0<x\leqslant3\end{cases}\]
$g(x)=[f(|x|)]+[|f(x)|]$ where $[.]$ denotes the greatest integer function. Which of the following statement(s) is/are true ?
\begin{qoptions}
\opt{A}{$g(1.7)=-1$}
\opt{B}{$g(x)$ is not differentiable at any point}
\opt{C}{$g(x)$ is continuous at all non integral points}
\opt{D}{$g(x)$ is differentiable at some points}
\end{qoptions}
\sect{SECTION-II(i)}
\paraq{1}{2}
A Rod of length 10cm slides between two lines, x-axis \& $y=-\sqrt3\,x$. If locus of mid point of rod is curve `S'.
\qn{1} x-intercept of S is
\qn{2} If maximum distance between origin and curve S is equal to $k$ then value of $k\sqrt3$ is
\paraq{3}{4}
Let $f(x)=\sin^{-1}(\sin e^x)$, $g(x)=\sin^{-1}\!\left(\dfrac{3^x+3^{-x}}{4}\right)$ and $h(x)=f(x)-g(x)$
\qn{3} If range of $g(x)$ is $\left[\dfrac{\pi}{a},\dfrac{\pi}{b}\right]$, then value of $(a+b)$ is
\qn{4} Number of solution(s) of $h(x)=0$ is
\paraq{5}{6}
Let $f:[0,\infty)\to[2,\infty)$ be a differentiable increasing and onto function satisfying
\[f^2(x)-2f(x)-\sqrt x=f^2(y)-2f(y)-\sqrt y,\ \forall x,y\in[0,\infty)\]
\qn{5} The value of $\displaystyle\lim_{x\to0^+}\dfrac{f(x)-2}{\sqrt x}$ is equal to -
\qn{6} If $g(x)$ be the inverse of $f(x)$, then $g'(4)$ is equal to -
\sect{SECTION-II(ii)}
\qn{7} A function $f$ satisfies $f(3x)=3f(x)\ \forall\ x\in[1,\infty)$ and $f(x)=1-|x-2|$ for $x\in[1,3]$. The smallest $x$ for which $f(x)=f(2001)$ is
\qn{8} For $\ell,m,n\in R$, $p$ and $q$ are the value of $\ell$ and $m$ respectively such that equation $\ell(13+11n)+m(19+17n)-4(8+7n)=0$ is satisfied $\forall\ n\in R$ then $(p+q)$ equals
\qn{9} The value of $k$ ($k>0$) such that the length of the longest interval in which the function $f(x)=\sin^{-1}|\sin kx|+\cos^{-1}(\cos kx)$ is constant is $\dfrac{\pi}{4}$ will be
\raggedcolumns
\newpage
\papersec{2}
\partsec{1}{PHYSICS}
\sect{SECTION-I}
\qn{1} In the given circuit, find the value of resistance `R' :
\qfig{p2_phy_s1_q01_1}
\begin{qoptions}
\opt{A}{20$\Omega$}
\opt{B}{25$\Omega$}
\opt{C}{15$\Omega$}
\opt{D}{10$\Omega$}
\end{qoptions}
\qn{2} In the given figure, $R=3\Omega$. Determine equivalent resistance between points A and B.
\qfig{p2_phy_s1_q02_1}
\begin{qoptions}
\opt{A}{3 $\Omega$}
\opt{B}{2 $\Omega$}
\opt{C}{1.5 $\Omega$}
\opt{D}{6 $\Omega$}
\end{qoptions}
\qn{3} A piece of conducting wire of resistance $R$ is cut into $2n$ equal parts. Half the parts are connected in series to form a bundle and remaining half in parallel to form another bundle. These bundles are then connected to give the maximum resistance. The resistance of the combination is
\begin{qoptions}
\opt{A}{$\dfrac R2\left(1+\dfrac1{n^2}\right)$}
\opt{B}{$\dfrac R2\left(1+n^2\right)$}
\opt{C}{$\dfrac{R}{2(1+n^2)}$}
\opt{D}{$R\left(n+\dfrac1n\right)$}
\end{qoptions}
\qn{4} In figure, points A and B are connected by a perfectly conducting wire. Calculate the current through AB
\qfig{p2_phy_s1_q04_1}
\begin{qoptions}
\opt{A}{2A}
\opt{B}{1 A}
\opt{C}{1.5 A}
\opt{D}{2.5 A}
\end{qoptions}
\qn{5} AB is a potentiometer wire of length 100 cm. When a cell $E_1$ is connected across AC, where AC = 75 cm, no current flows from $E_2$. The internal resistance of the cell $E_1$ is negligible.
\qfig{p2_phy_s1_q05_1}
Potential gradient along AB and value of $E_2$ respectively is.
\begin{qoptions}
\opt{A}{0.01 V cm$^{-1}$, 1.2 volt}
\opt{B}{0.03 Vcm$^{-1}$, 1.3 volt}
\opt{C}{0.04 Vm$^{-1}$, 1.0 volt}
\opt{D}{0.02 Vcm$^{-1}$, 1.5 volt}
\end{qoptions}
\qn{6} Find the reading of the ideal ammeter connected in the given circuit. Assume that the cells have negligible internal resistance.
\qfig{p2_phy_s1_q06_1}
\begin{qoptions}
\opt{A}{2.5 A}
\opt{B}{5 A}
\opt{C}{2 A}
\opt{D}{none of these}
\end{qoptions}
\qn{7} A solid sphere of uniform density \& having total mass $(M)$ now two spherical cavity is removed as shown in figure. Radius of sphere is 4 meter. Radius of each spherical cavity is 1 meter. Value of gravitation field at centre of sphere will be :-
\qfig{p2_phy_s1_q07_1}
\begin{qoptions}
\opt{A}{$\dfrac{31GM}{1024}$}
\opt{B}{$\dfrac{31GM}{256}$}
\opt{C}{$\dfrac{GM}{1024}$}
\opt{D}{Zero}
\end{qoptions}
\qn{8} Four similar masses orbiting around a circle of radius `r'. Due to their mutual gravitation attraction force velocity of each particle will be :-
\qfig{p2_phy_s1_q08_1}
\begin{qoptions}
\opt{A}{$\sqrt{\dfrac{Gm}{r}\left(\dfrac{1+2\sqrt2}{4}\right)}$}
\opt{B}{$\sqrt{\dfrac{Gm}{r}\left(1+2\sqrt2\right)}$}
\opt{C}{$\sqrt{\dfrac{Gm}{r}}$}
\opt{D}{$\sqrt{\dfrac{Gm}{r}\left(\dfrac{1+2\sqrt2}{2}\right)}$}
\end{qoptions}
\qn{9} Find out flux passing through shaded surface of the sphere when a point charge `q' is placed at centre of sphere. [Radius of sphere is $R$]
\qfig{p2_phy_s1_q09_1}
\begin{qoptions}
\opt{A}{$\dfrac{q}{\varepsilon_0}$}
\opt{B}{$\dfrac{q}{4\varepsilon_0}$}
\opt{C}{$\dfrac{q}{2\varepsilon_0}$}
\opt{D}{Zero.}
\end{qoptions}
\qn{10} A cube of side `$\ell$' placed as shown in figure $E_x=\beta x^{\frac12}$. Flux passing through cube will be:-
\qfig{p2_phy_s1_q10_1}
\begin{qoptions}
\opt{A}{$2\sqrt2\,\beta\ell^{\frac52}$}
\opt{B}{$(\sqrt2-1)\,\beta\ell^{\frac52}$}
\opt{C}{$-\beta\ell^{\frac52}$}
\opt{D}{Zero}
\end{qoptions}
\qn{11} Two solid spheres having volume charge density $\rho$ and $-\rho$ respectively as shown in figure. Find out the potential difference $(\Delta V)$ between centre of both sphere for $d=R$
\qfig{p2_phy_s1_q11_1}
\begin{qoptions}
\opt{A}{$\dfrac{\rho}{3\varepsilon_0}R^2$}
\opt{B}{$\dfrac{\rho}{\varepsilon_0}R^2$}
\opt{C}{0}
\opt{D}{$\dfrac{2\rho}{\varepsilon_0}R^2$}
\end{qoptions}
\qn{12} There is a infinite chain of charge $q$ \& $-q$. Distance between the two neighbouring charge is equal to `d'. The interaction energy of one of charge with all other charge is :
\qfig{p2_phy_s1_q12_1}
\begin{qoptions}
\opt{A}{$\dfrac{-2q^2}{4\pi\varepsilon_0d}$}
\opt{B}{$\dfrac{-2q^2\ell n2}{4\pi\varepsilon_0d}$}
\opt{C}{$\dfrac{2q^2\ell n2}{4\pi\varepsilon_0d}$}
\opt{D}{$\dfrac{q^2}{\pi\varepsilon_0d}$}
\end{qoptions}
\qn{13} Three arc are joined together as shown in figure having charges $Q$, $-2Q$, $3Q$ respectively. Find out the electric potential at point C.
\qfig{p2_phy_s1_q13_1}
\begin{qoptions}
\opt{A}{$\dfrac{Q}{2\pi\varepsilon_0R}$}
\opt{B}{$\dfrac{5Q}{12\pi\varepsilon_0R}$}
\opt{C}{$\dfrac{3Q}{32\pi\varepsilon_0R}$}
\opt{D}{$\dfrac{3Q}{8\pi\varepsilon_0R}$}
\end{qoptions}
\qn{14} Instantaneous speed on end A of rod is `v'. Angular velocity of rod of length `L' will be.
\qfig{p2_phy_s1_q14_1}
\begin{qoptions}
\opt{A}{$\dfrac{v}{2L}$}
\opt{B}{$\dfrac{v\sqrt3}{2L}$}
\opt{C}{$\dfrac{v}{L}$}
\opt{D}{$\dfrac{4v}{2L}$}
\end{qoptions}
\qn{15} Equilateral prism of mass $m$ is rest on rough horizontal floor. If coefficient of friction is sufficiently high such that prism does not slide before topple. Minimum force required to topple prism is (coefficient is friction = $\mu$)
\qfig{p2_phy_s1_q15_A}
\begin{qoptions}
\opt{A}{$\dfrac{mg}{\sqrt3}$}
\opt{B}{$\dfrac{\mu mg}{\sqrt3}$}
\opt{C}{$\dfrac{mg}{4}$}
\opt{D}{$\dfrac{\mu mg}{4}$}
\end{qoptions}
\qn{16} A force `F' act tangentially at highest point of a disc of mass `m' \& radius `R'. If disc is rolling without sliding. Acceleration of centre of disc will be
\qfig{p2_phy_s1_q16_1}
\begin{qoptions}
\opt{A}{$\dfrac{2F}{3m}$}
\opt{B}{0}
\opt{C}{$\dfrac{10F}{7m}$}
\opt{D}{$\dfrac{4F}{3m}$}
\end{qoptions}
\qn{17} Time period of a simple pendulum whose length is $\dfrac R2$, where $R$ is radius of the earth oscillating near earth's surface will be -
\begin{qoptions}
\opt{A}{$2\pi\sqrt{\dfrac{R}{2g}}$}
\opt{B}{$2\pi\sqrt{\dfrac{R}{g}}$}
\opt{C}{$2\pi\sqrt{\dfrac{R}{3g}}$}
\opt{D}{$2\pi\sqrt{\dfrac{3R}{g}}$}
\end{qoptions}
\qn{18} A point charge $-q$ is placed inside cavity of a metallic sphere of radius $R$. Radius of the cavity is $r$. Potential at point P will be -
\qfig{p2_phy_s1_q18_1}
\begin{qoptions}
\opt{A}{$\dfrac{-q}{4\pi\varepsilon_0y}$}
\opt{B}{$\dfrac{-q}{4\pi\varepsilon_0y}+\dfrac{q}{4\pi\varepsilon_0x}$}
\opt{C}{$\dfrac{-q}{4\pi\varepsilon_0x}$}
\opt{D}{$\dfrac{q}{4\pi\varepsilon_0y}-\dfrac{q}{4\pi\varepsilon_0x}$}
\end{qoptions}
\qn{19} A particle is projected from bottom of a blind narrow tunnel of depth $\dfrac R4$ made along a radius inside the earth. If minimum velocity needed to escape out of the gravitational field of the earth is $\dfrac14\sqrt{\dfrac{xGM}{R}}$. Find $x$. $M$ and $R$ are mass and radius of the earth respectively.
\begin{qoptions}
\opt{A}{39}
\opt{B}{78}
\opt{C}{86}
\opt{D}{0}
\end{qoptions}
\qn{20} A ray of light incidents at an angle of incidence $\theta$ in water at water-air interface. Graph between angle of deviation $(\delta)$ and angle of incidence $(\theta)$ will be-
\qfig{p2_phy_s1_q20_1}
\begin{qoptions}
\opt{A}{\qfigi{p2_phy_s1_q20_A}}
\opt{B}{\qfigi{p2_phy_s1_q20_B}}
\opt{C}{\qfigi{p2_phy_s1_q20_C}}
\opt{D}{\qfigi{p2_phy_s1_q20_D}}
\end{qoptions}
\sect{SECTION-II}
\qn{1} A body is thrown from surface of earth with velocity $\left(\sqrt{\dfrac{gR_e}{2}}\right)$ [$R_e$ radius of earth] at some angle from vertical. Maximum height reach by the body is $\dfrac{R_e}{4}$. If angle of projection is $\cos^{-1}\!\left[\dfrac{\sqrt x}{4}\right]$ find the value of x.
\qn{2} A uniform solid sphere having mass density is $\rho$. The ratio of gravitational pressure at a distance $r=R/2$ to the centre of sphere is $\dfrac{3x}{8}$, the value of x will be
\qfig{p2_phy_s2_q02_1}
\qn{3} Linear charge density of ring of radius `R' given as $\lambda=\lambda_0\cos\dfrac\theta2$. The potential at centre of ring will be
\qfig{p2_phy_s2_q03_1}
\qn{4} Uniform $\vec E=100$v/m is directed as shown in figure. The potential difference between $v_A-v_O$ is
\qfig{p2_phy_s2_q04_1}
\qn{5} A ray incident on medium boundary as shown in figure. Unit vector in direction of incident ray is given as $\dfrac1{\sqrt2}\hat i+\dfrac1{\sqrt2}\hat j$ \& unit vector of ray in medium (2) is given by $\dfrac12\hat i+\dfrac{\sqrt3}2\hat j$. The ratio of
\[\dfrac{\left(\frac{\mu_1}{\mu_2}\right)^2+\frac12}{\left(\frac{\mu_1}{\mu_2}\right)^2-\frac12}\]
is
\qfig{p2_phy_s2_q05_1}
\qn{6} A ray incident normal to the surface of prism as shown in figure. Angle of deviation of ray from the mirror w.r.t incident ray is $210^\circ-\sin^{-1}\dfrac x4$, find the value of x.
\qfig{p2_phy_s2_q06_1}
\qn{7} Three identical rod each of mass $m$ \& length $\ell$ form a equilateral triangle. Moment of inertia about given axis is $\dfrac{m\ell^2}{x}$, value of x will be :-
\qfig{p2_phy_s2_q07_1}
\qn{8} In given diagram a uniform plank weight is 500 N. Find out weight W if $T_1$ \& $T_2$ are equal.
\qfig{p2_phy_s2_q08_1}
\qn{9} Figure shows segment of a current carrying wire made of copper in steady state. If ratio of magnitudes of electric field at cross-section $S_1$ and $S_2$ is $5n$. Find value of n.
\qfig{p2_phy_s2_q09_1}
\qn{10} A thick conducting shell of inner radius $R$ and outer radius $2R$ has charge $Q$ at its center. It is covered by a thin conducting shell of radius $4R$ which has charge $2Q$. If charge flown into the earth when switch is closed is $nQ$, find $n$.
\qfig{p2_phy_s2_q10_1}
\raggedcolumns
\partsec{2}{CHEMISTRY}
\sect{SECTION-I}
\qn{1} The correct order of second I.E. of C, N, O and F are in the order
\begin{qoptions}
\opt{A}{F $>$ O $>$ N $>$ C}
\opt{B}{C $>$ N $>$ O $>$ F}
\opt{C}{O $>$ N $>$ F $>$ C}
\opt{D}{O $>$ F $>$ N $>$ C}
\end{qoptions}
\qn{2} Calamine malachite, magnetite and cryolite, respectively are:
\begin{qoptions}
\opt{A}{ZnSO$_4$, CuCO$_3$, Fe$_2$O$_3$, AlF$_3$}
\opt{B}{ZnSO$_4$, Cu(OH)$_2$, Fe$_3$O$_4$, Na$_3$AlF$_6$}
\opt{C}{ZnCO$_3$, CuCO$_3\cdot$Cu(OH)$_2$, Fe$_3$O$_4$, Na$_3$AlF$_6$}
\opt{D}{ZnCO$_3$, CuCO$_3$, Fe$_2$O$_3$, Na$_3$AlF$_6$}
\end{qoptions}
\qn{3} Select the incorrect IUPAC name from the following
\begin{qoptions}
\opt{A}{H[AuCl$_4$] : Hydrogen tetrachloridoaurate(III)}
\opt{B}{[PtCl$_6$]$^{2-}$ : Hexachloridoplatinate(IV) ion}
\opt{C}{[Cu(NH$_3$)$_4$]SO$_4$ : Tetraamminecopper(II) sulphate}
\opt{D}{[Cr(en)$_3$]Cl$_3$ : Tris (ethylenediamine)chromium(III) chloride}
\end{qoptions}
\qn{4} Spin only magnetic moment in BM of [Fe(CO)$_4$(C$_2$O$_4$)]$^+$ is :
\begin{qoptions}
\opt{A}{5.92}
\opt{B}{0}
\opt{C}{2.82}
\opt{D}{1.73}
\end{qoptions}
\qn{5} Which of the following complex is optically active :
\begin{qoptions}
\opt{A}{trans-[Co(NH$_3$)$_4$Cl$_2$]$^+$}
\opt{B}{[Cr(H$_2$O)$_6$]$^{3+}$}
\opt{C}{cis-[Co(NH$_3$)$_2$(en)$_2$]$^{3+}$}
\opt{D}{trans-[Co(NH$_3$)$_2$(en)$_2$]$^{3+}$}
\end{qoptions}
\qn{6} Select INCORRECT order against property indicated--
\begin{qoptions}
\opt{A}{AsCl$_3<$ SbCl$_3<$ BiCl$_3$ (relative stability)}
\opt{B}{ZnO $>$ MgO $>$ CaO $>$ BaO (Covalent character)}
\opt{C}{NaF $<$ MgF$_2<$ AlF$_3$ (M.P.) Melting point}
\opt{D}{BeF$_2>$ MgF$_2>$ CaF$_2>$ SrF$_2>$ BaF$_2$ (water solubility)}
\end{qoptions}
\qn{7} Which of the following gives alcohol as product :
\begin{qoptions}
\opt{A}{\qfigi{p2_chem_s1_q07_A}}
\opt{B}{CH--CN $\xrightarrow[\mathrm{HOH}]{\mathrm{SnCl_2,\ HCl}}$}
\opt{C}{\qfigi{p2_chem_s1_q07_C}}
\opt{D}{\qfigi{p2_chem_s1_q07_D}}
\end{qoptions}
\qn{8} Select reaction from the following in which only CO$_2$ gas is released :
\begin{qoptions}
\opt{A}{\qfigi{p2_chem_s1_q08_A}}
\opt{B}{\qfigi{p2_chem_s1_q08_B}}
\opt{C}{\qfigi{p2_chem_s1_q08_C}}
\opt{D}{\qfigi{p2_chem_s1_q08_4}}
\end{qoptions}
\qn{9} \qfig{p2_chem_s1_q09_1}
Major product of above reaction :
\begin{qoptions}
\opt{A}{\qfigi{p2_chem_s1_q09_A}}
\opt{B}{CH$_3$--CH$_3$}
\opt{C}{CH$_2$=CH$_2$}
\opt{D}{HC$\equiv$CH}
\end{qoptions}
\qn{10} Following conversion can not be carried out by using sequence
\[\mathrm{CH_3{-}CH_2{-}\overset{O}{\overset{\|}{C}}{-}OH\ \longrightarrow\ CH_3{-}CH_2{-}\overset{O}{\overset{\|}{C}}{-}H}\]
\begin{qoptions}
\opt{A}{(i) LiAlH$_4$ , PBr$_3$ , DMSO /NaHCO$_3$ / $\Delta$}
\opt{B}{(i) SOCl$_2$ (ii) CH$_3$MgX (excess)}
\opt{C}{(i) Ca(OH)$_2$ (ii) (HCOO)$_2$Ca/ Dry distillation}
\opt{D}{(i) LiAlH$_4$ (ii) Cu/300$^\circ$C}
\end{qoptions}
\qn{11} Which of the following will not undergo periodic cleavage with HIO$_4$, $\Delta$?
\begin{qoptions}
\opt{A}{\qfigi{p2_chem_s1_q11_A}}
\opt{B}{\qfigi{p2_chem_s1_q11_B}}
\opt{C}{\qfigi{p2_chem_s1_q11_C}}
\opt{D}{\qfigi{p2_chem_s1_q11_D}}
\end{qoptions}
\qn{12} \qfig{p2_chem_s1_q12_1}
Product P is?
\begin{qoptions}
\opt{A}{\qfigi{p2_chem_s1_q12_A}}
\opt{B}{\qfigi{p2_chem_s1_q12_B}}
\opt{C}{\qfigi{p2_chem_s1_q12_C}}
\opt{D}{\qfigi{p2_chem_s1_q12_D}}
\end{qoptions}
\qn{13} Which of the following ether can not be synthesised by Williamson ether synthesis in good yield :
\begin{qoptions}
\opt{A}{Ph--O--CH$_2$--CH$_3$}
\opt{B}{\qfigi{p2_chem_s1_q13_B}}
\opt{C}{CH$_3$--O--CH$_3$}
\opt{D}{Ph--CH$_2$--O--Ph}
\end{qoptions}
\qn{14} At 300 K, the vapour pressure of an ideal solution containing 3 mole of A and 2 mole of B is 600 torr. At the same temperature, if 1.5 mole of A \& 0.5 mole of C (non-volatile) are added to this solution the vapour pressure of solution increases by 30 torr. What is the value of $P_B^{\circ}$?
\begin{qoptions}
\opt{A}{940 torr}
\opt{B}{405 torr}
\opt{C}{90 torr}
\opt{D}{100 torr}
\end{qoptions}
\qn{15} What is the correct sequence of osmotic pressures of 0.01 M aq. solutions of :\par
(i) Al$_2$(SO$_4$)$_3$ ($\pi_1$)\par
(ii) Na$_3$PO$_4$ ($\pi_2$)\par
(iii) BaCl$_2$ ($\pi_3$)\par
(iv) Glucose ($\pi_4$)\par
[Assume 100\% ionisation of electrolytes]
\begin{qoptions}
\opt{A}{$\pi_4>\pi_2>\pi_3>\pi_1$}
\opt{B}{$\pi_3>\pi_4>\pi_2>\pi_1$}
\opt{C}{$\pi_3>\pi_4>\pi_1>\pi_2$}
\opt{D}{$\pi_1>\pi_2>\pi_3>\pi_4$}
\end{qoptions}
\qn{16} Find ratio of mole fractions of gases N$_2$ \& O$_2$ dissolved in water from air at 25$^{\circ}$ \& 1 atm. Assume air is 20\,\% by volume of O$_2$ \& 80\% by volume of N$_2$\par
\textbf{Given :} $K_H(\mathrm{N_2})=2\times10^4$ atm\par
\hspace*{3.3em}$K_H(\mathrm{O_2})=10^4$ atm
\begin{qoptions}
\opt{A}{8 : 1}
\opt{B}{1 : 8}
\opt{C}{2 : 1}
\opt{D}{1 : 1}
\end{qoptions}
\qn{17} For a process A + B $\to$ P, the reaction is second order with respect to A and zero order with respect to B. When 1 mole each of A \& B are taken in 1 litre vessel the initial rate is $1\times10^{-2}$ mol/L-sec. The rate of reaction when 50\% of the reactant have been converted to product would be -
\begin{qoptions}
\opt{A}{$1\times10^{-2}$ mol/L-sec}
\opt{B}{$2.5\times10^{-3}$ mol/L-sec}
\opt{C}{$5\times10^{-2}$ mol/L-sec}
\opt{D}{$5\times10^{-3}$ mol/L-sec}
\end{qoptions}
\qn{18} A vessel contains He gas at 20 atm and CH$_4$ gas at `P' atm. When a pin-hole is made in the vessel, the gases coming out initially contains 80\% He gas, by mole. The value of `P' is :-
\begin{qoptions}
\opt{A}{1 atm}
\opt{B}{40 atm}
\opt{C}{2.5 atm}
\opt{D}{10 atm}
\end{qoptions}
\qn{19} A real gas behaves like an ideal gas at 132$^{\circ}$C in low pressure range. Critical temperature for this gas is -
\begin{qoptions}
\opt{A}{120 K}
\opt{B}{39 K}
\opt{C}{132 K}
\opt{D}{15 K}
\end{qoptions}
\qn{20} In a CCP lattice of X \& Y, X atoms are present at the corners while Y atoms are at face centers. Then the formula of the compound would be if one of the X-atoms from a corner in each unit cell is replaced by Z-atom :
\begin{qoptions}
\opt{A}{X$_7$ Y$_{24}$ Z$_2$}
\opt{B}{X$_7$ Y$_{24}$ Z}
\opt{C}{X$_{24}$ Y$_7$ Z}
\opt{D}{X Y$_{24}$ Z}
\end{qoptions}
\sect{SECTION-II}
\qn{1} Find out the no. of correct statements :\par
(a) Bond length N$_2^+>$ Bond length N$_2$\par
(b) Bond length NO$^+<$ Bond length of NO\par
(c) Bond length CN$^-<$ Bond length of CN\par
(d) Bond length O$_2^-<$ Bond length of O$_2^{-2}$\par
(e) Bond length O$_2>$ Bond length of O$_2^+$\par
(f) Bond length B$_2>$ Bond length of B$_2^-$
\qn{2} The number of 2c -- 2e bonds present in a dimer molecule of Al(CH$_3$)$_3$ is-
\qn{3} Find the total number of type of isomerism (from given below) shown by the complex\par
\textbf{[Fe(NH$_3$)$_4$(H$_2$O)Br] (NO$_3$)$_2$}\par
Ionization isomerism, linkage isomerism, ligand isomerism, Hydrate isomerism, geometrical isomerism
\qn{4} Total number of geometrical isomers for the complex [RhCl(CO)(PPh$_3$) (NH$_3$)] is:
\qn{5} Calculate total number of moles of Grignard (MeMgCl) required for complete reaction with following compound ?
\qfig{p2_chem_s2_q05_1}
\qn{6} How many of the following will give positive Fehling test :
\qfig{p2_chem_s2_q06_1}
\qn{7} How many of the following will not give iodoform test :
\qfig{p2_chem_s2_q07_1}
\qn{8} For a reaction A$\to$P\par
Half life measured for two different values of initial concentrations of A, $5\times10^{-3}$ M and $25\times10^{-4}$ M are 1 and 8 hrs respectively. If initial concentration is adjusted to $1.25\times10^{-3}$ M, what is the half-life? (in hrs)
\qn{9} In the following reaction : $x\mathrm{A}\longrightarrow y\mathrm{B}$, following equation is given :
\[\log\left[-\frac{d[\mathrm{A}]}{dt}\right]=\log\left[\frac{d[\mathrm{B}]}{dt}\right]+0.3,\]
Find ratio $\left(\dfrac yx\right)$ : [Given $\log2=0.3$]
\qn{10} A metal `A' exists as BCC structure having uncovered length between the adjacent atoms along the edge in an unit cell equal to 1.2 \AA{}. If molar mass of the metal is 307.2 gm/mole then calculate density of the crystal in gm/ml. [Use : $N_A=6\times10^{23}$ and $\sqrt3=1.7$]
\raggedcolumns
\partsec{3}{MATHEMATICS}
\sect{SECTION-I}
\qn{1} $\displaystyle\lim_{n\to\infty}\dfrac{\displaystyle\sum_{r=0}^{n}\tan^{-1}\left(1+r+r^2\right)}{n}$ is equal to-
\begin{qoptions}
\opt{A}{1}
\opt{B}{2}
\opt{C}{$\dfrac\pi4$}
\opt{D}{$\dfrac\pi2$}
\end{qoptions}
\qn{2} The complete set of values the parameter `a' so that the point $P\left(a,\ \dfrac{1}{1+a^2}\right)$ doesn't lie outside the triangle formed by the lines $15y=x+1$, $78y=118-23x$ and $y+2=0$ is -
\begin{qoptions}
\opt{A}{$(0,5)$}
\opt{B}{$[2,5]$}
\opt{C}{$(1,5)$}
\opt{D}{$[0,2]$}
\end{qoptions}
\qn{3} Which of the following is true-
\begin{qoptions}
\opt{A}{If number of elements in set A are finite, such that $f:A\to A$ is an one-one function, then $f$ is necessarily onto.}
\opt{B}{If a function is continuous in its domain and changes its sign for any 2 values of $x$ then odd number of roots lie between the given values of $x$}
\opt{C}{If $f:A\to A$ is one-one then it must be onto}
\opt{D}{It is possible that a curve contains a local maxima and a global minima at the same point.}
\end{qoptions}
\qn{4} $\displaystyle\int\dfrac{2\left(x^3-1\right)}{x\left(2x^3+1\right)}dx$ is equal to- (where C is the constant of integration)
\begin{qoptions}
\opt{A}{$\ell n\left(2x+\dfrac1{x^2}\right)+C$}
\opt{B}{$\ell n\left(2x^3+1\right)+C$}
\opt{C}{$\tan^{-1}\dfrac{1}{x^{3/2}+1}+\ell n\left(\dfrac{x-1}{x+1}\right)+C$}
\opt{D}{$\ell n\left(x+\sqrt{\dfrac{x-1}{x+1}}\right)+C$}
\end{qoptions}
\qn{5} If the function $g(x)=\begin{cases}ae^x, & x\leqslant0\\ b\cos x+x, & x>0\end{cases}$ is differentiable, then the value of $a^2+b^2$ is -
\begin{qoptions}
\opt{A}{5}
\opt{B}{2}
\opt{C}{1}
\opt{D}{13}
\end{qoptions}
\qn{6} The number of real solutions of the equation $\tan(e^x)=e^x+e^{-x}$, $x>0$, is
\begin{qoptions}
\opt{A}{0}
\opt{B}{1}
\opt{C}{2}
\opt{D}{none of these}
\end{qoptions}
\qn{7} If the normal to the curve $y=f(x)$ at the point $(4,6)$ makes an angle $\dfrac{2\pi}{3}$ with the positive $x$ axis in anticlockwise direction, then $f'(4)$ is
\begin{qoptions}
\opt{A}{$\sqrt3$}
\opt{B}{$-\sqrt3$}
\opt{C}{$\dfrac{1}{\sqrt3}$}
\opt{D}{$\dfrac{-1}{\sqrt3}$}
\end{qoptions}
\qn{8} $\displaystyle\lim_{x\to0}\dfrac{5\sin x+x\cos x}{2\tan x-x^2}$ is
\begin{qoptions}
\opt{A}{$\dfrac52$}
\opt{B}{3}
\opt{C}{0}
\opt{D}{does not exist}
\end{qoptions}
\qn{9} If linear functions $f(x)$ and $g(x)$ satisfy $\displaystyle\int\left[(1-2x)\cos x+(3+2x)\sin x\right]dx=f(x)\sin x+g(x)\cos x+C$ (where C is constant of integration), then
\begin{qoptions}
\opt{A}{$f(x)=3-2x$}
\opt{B}{$g(x)=5-2x$}
\opt{C}{$f(x)=3+2x$}
\opt{D}{$g(x)=-5+2x$}
\end{qoptions}
\qn{10} Let $f:R\to R$, $f(x)=ax^3+bx^2+cx+d$, has no extreme value, then which of the following is always correct ?
\begin{qoptions}
\opt{A}{$3a+2b+c\geq0$}
\opt{B}{$c\leq0$}
\opt{C}{$(3a+2b+c)c\geq0$}
\opt{D}{$3a+2b+2c\geq0$}
\end{qoptions}
\qn{11} If $f(x)=3-x^2$, $1\leq x\leq4$, then the domain of $\log_e(f(2x))$ is
\begin{qoptions}
\opt{A}{$\left(\dfrac{-\sqrt3}{2},\dfrac{\sqrt3}{2}\right)$}
\opt{B}{$\left[\dfrac12,2\right]$}
\opt{C}{$\left(0,\dfrac{\sqrt3}{2}\right)$}
\opt{D}{$\left[\dfrac12,\dfrac{\sqrt3}{2}\right)$}
\end{qoptions}
\qn{12} If two adjacent sides of a parallelogram are given by $2x^2+2y^2-5xy=0$ and $5x+2y-1=0$ is one of its diagonal, then area of parallelogram is $\dfrac1k$ sq. units, where $k$ is
\begin{qoptions}
\opt{A}{16}
\opt{B}{72}
\opt{C}{18}
\opt{D}{36}
\end{qoptions}
\qn{13} Number of positive integral values of `$k$' for which the equation $k=|x+|2x-1||-|x-|2x-1||$ has exactly three real solutions, is
\begin{qoptions}
\opt{A}{0}
\opt{B}{2}
\opt{C}{3}
\opt{D}{5}
\end{qoptions}
\qn{14} $\displaystyle\int\left(6x^2+5x+4\right)\left(x^2+x+1\right)^6\cdot x^{27}dx$ equals (where C is integration constant)
\begin{qoptions}
\opt{A}{$\dfrac{x^{28}\left(1+x+x^2\right)^7}{7}+C$}
\opt{B}{$\dfrac{x^{24}\left(1+x+x^2\right)^7}{7}+C$}
\opt{C}{$\dfrac{x^{24}\left(1+x+x^2\right)^6}{7}+C$}
\opt{D}{$\dfrac{x^{28}\left(1+x+x^2\right)^8}{7}+C$}
\end{qoptions}
\qn{15} The value of $\displaystyle\lim_{x\to\infty}x^{\frac13}\left((x+1)^{\frac23}-(x-1)^{\frac23}\right)$ is
\begin{qoptions}
\opt{A}{$\dfrac43$}
\opt{B}{$\dfrac{-1}{3}$}
\opt{C}{$\dfrac13$}
\opt{D}{$\dfrac{-2}{3}$}
\end{qoptions}
\qn{16} In the interval $(0,1)$ maximum value of the function $f(x)=|x\,\ell nx|$ is
\begin{qoptions}
\opt{A}{0}
\opt{B}{1}
\opt{C}{$e$}
\opt{D}{$e^{-1}$}
\end{qoptions}
\qn{17} If the equation $\sin^{-1}(x-1)+\cos^{-1}(x-3)+\tan^{-1}\left(\dfrac{x}{-x^2+2}\right)=m$ holds, then the value of `$m$' is
\begin{qoptions}
\opt{A}{$\dfrac{-3\pi}{4}$}
\opt{B}{$\dfrac{3\pi}{4}$}
\opt{C}{$\dfrac{-5\pi}{4}$}
\opt{D}{$\dfrac{5\pi}{4}$}
\end{qoptions}
\qn{18} The expression $x^3-3x^2-9x+c$ can be written in the form $(x-a)^2(x-b)$ if the values of c is
\begin{qoptions}
\opt{A}{$-5$ or 27}
\opt{B}{5 or $-27$}
\opt{C}{5 or 27}
\opt{D}{$-5$ or $-27$}
\end{qoptions}
\qn{19} A real valued function $y=f(x)$ satisfies the relation $f\left(x-\dfrac49\right)+2x\leqslant\dfrac94x^2+\dfrac89\leqslant f\left(x+\dfrac49\right)-2x$. The value of $f''(2)$ is
\begin{qoptions}
\opt{A}{4}
\opt{B}{$\dfrac92$}
\opt{C}{$\dfrac{15}2$}
\opt{D}{$\dfrac{27}2$}
\end{qoptions}
\qn{20} Number of points where the function $f(x)=\max\left(\sqrt{2x-x^2},\ 2-x\right)$ is non differentiable, is
\begin{qoptions}
\opt{A}{0}
\opt{B}{1}
\opt{C}{2}
\opt{D}{3}
\end{qoptions}
\sect{SECTION-II}
\qn{1} If $m$ is the least value of $f(x)=\sqrt{x^2-2x+2}+\sqrt{x^2-4x+29}$, then $[m]$ is equal to\par
(where $[.]$ denotes greatest integer function)
\qn{2} If `a' and `b' are the greatest and least values of $f(x)=18x^{4/3}-9x^{1/3}+1$ in the interval $[-1,8]$ then the value of $(a+112b)$ is
\qn{3} If $f(x)=10x-2x^2$, then
\[\lim_{h\to0}\frac{f(3\cosh h+4\sinh h-2)-8}{f(3e^h-5\sec h+4)-12}\]
is equal to
\qn{4} If $f(x)=\begin{cases}\dfrac{e^{ax}+\ell n(1+4x)+b}{x^3} & x\neq0\\ L & x=0\end{cases}$ is continuous at $x=0$, then $12abL$ is equal to
\qn{5} If $y=\left(\ell n\left(x+\sqrt{x^2+1}\right)\right)^2$ and
\[(x^2+1)\frac{d^3y}{dx^3}+(px+q)\frac{d^2y}{dx^2}+r\frac{dy}{dx}=sy,\]
then value of $(p+q+r+s)$ is
\qn{6} $f(x)=\begin{cases}(x+1)^2 & -3<x\leqslant0\\ 2x^{1/3}+a & 0<x\leqslant1\\ 2-x & 1<x<2\end{cases}$ has maxima at $x=0,1$. If number of integral values of `a' is $L$ and total number of extrema of $f(x)$ is $M$ then $(L+M)$ is
\qn{7} If $f(x),g(x),h(x)$ be such that $f'(x)=g(x)$, $g'(x)=h(x)$, $h'(x)=f(x)$, $f(0)=4$, $g(0)=5$, $h(0)=0$ then the value of $\left(f^3(x)+g^3(x)+h^3(x)-3f(x)g(x)h(x)\right)$ at $x=2012$ is
\qn{8} Let $y=f(x)$ be a curve passing through $P(0,1)$ \& $Q\left(\dfrac12,\dfrac{17}{16}\right)$ and having a property that slope of the tangent at any point is proportional to the cube of the abscissa at the point. If length of normal at Q be $\ell$ then $\dfrac{32\ell}{\sqrt5}$ is equal to
\qn{9} Number of integers in the range of function $f(x)=\dfrac{100}{3\left(x^6+12x^4+5x^2+1\right)}$ is
\qn{10} If $\displaystyle\int\dfrac{\sqrt{e^{7x}}}{\sqrt[4]{e^{2x}+e^{-2x}}}dx=\dfrac1k\left(e^{\alpha x}+1\right)^{\beta/\gamma}+C$ (where C is constant of integration \& $k$, $\alpha$, $\beta$, $\gamma$ are positive integers and $\beta$, $\gamma$ are relatively prime) then $\left(\dfrac k\beta+\dfrac\gamma\alpha\right)$ is
\end{multicols}
\end{document}
